\subsection{Phase 3 - Rhetorical role prediction}
The initial role-prediction runs compare LoRA, rank-stabilized LoRA, IA$^3$, LoRMA, prompts and distillation. A later experiment fully tunes bases~03/01 with three-layer bidirectional GRU document heads, then selects a blend with earlier models on dev data. The separate HighF1 branch uses target-focused 448-token inputs, attention/mean/max pooling, position and domain embeddings, and a two-layer bidirectional GRU over twelve-sentence windows. Pair heads distinguish argument parties, precedent reliance and the source of a court ruling; a CRF adds transition scores. Five document folds favored raw decoding. Corporate Law and IT share the model, with a domain feature to distinguish them.

\subsection{Phase 4 - Court judgment prediction}
For judgment prediction, each case is represented by up to eight 640-token chunks. Features are projected to width 256, combined with position embeddings and passed through a one-layer chunk Transformer. Attention, mean and max pooling feed a binary classifier, with its threshold selected by macro F1. The six variants use full base~01, partial base~03, LoRMA base~04, adaptive-rank base~02, adapters base~03 and distilled base~03. Partial base~03 gives the best saved dev result at epoch~3 and threshold~0.55. The intended grouped outer folds did not finish, so this is a dev-selected result.

\subsection{Phase 5 - Legal statute identification}
Statute prediction combines pooled document features with the strongest chunk evidence for each label. Training uses four 512-token chunks; inference allows twelve, selected for coverage and statute cues. Sigmoids and calibrated thresholds return a label set. The six variants are partial base~03, full base~01, LoRMA base~04, adapters base~02, adaptive-rank base~03 and distilled base~03. Repeated gold statute IDs are deduplicated. All six reports are complete: slot~01 leads holdout macro F1 at 0.37934, with micro F1 0.48132 and subset accuracy 0.04143. Distillation in slot~06 is close at 0.37650 macro F1 and 0.47888 micro F1, with higher subset accuracy, 0.04515, and 32.0\% fewer trainable parameters, 15.08M versus 22.18M.

\subsection{Phase 6 - Abstractive summarization}
The summarization runs use LoRA, adaptive rank, LoRMA, Houlsby adapters, Pfeiffer adapters and sequence distillation on the original foundations. Source selection keeps section and legal-cue evidence within each model's context limit. References are capped at 480 tokens. Decoder-only models generate after the packed source; base~02 uses cross-attention. Greedy cached decoding includes length hints, a minimum output budget and repeated-ngram suppression. There is no separate candidate reranker. For distillation, a deterministic subset of fitting targets is replaced with summaries from the project teacher. BERTScore~\cite{zhang2020bertscore} is used only for evaluation.

\subsection{Phase 7 - Prior Case Retrieval (PCR)}
Retrieval uses a shared encoder to turn up to six 384-token chunks into normalized 256-dimensional vectors. The neural score gives equal weight to document cosine and the average best candidate-chunk match for each query chunk. Reciprocal rank fusion combines this ranking with full-document and passage BM25~\cite{cormack2009rrf}. The six variants use full base~03, partial base~01, LoRMA base~04, IA$^3$ base~02, prompted base~03 and distilled base~02. Known positives are excluded from negatives; one mining refresh and a fitting-only teacher provide additional supervision. Every output contains seven unique IDs. The proposed cross-encoder branch was not implemented.
